% !TeX program = pdflatex
% Official Elsevier elsarticle class and numerical template.
% build.py also generates a single-column [preprint,11pt,sort&compress] copy.
\documentclass[final,5p,times,twocolumn,sort&compress]{elsarticle}
\usepackage[T1]{fontenc}
\usepackage{amsmath,amssymb,amsthm,booktabs,array,longtable}
\usepackage{algorithm,algpseudocode,placeins,xurl}
\usepackage{xr-hyper}
\IfFileExists{supplement_labels.tex}{% Generated cross-document labels; rebuild with build.py after edits.
\newlabel{S-sec:extended_selection}{{S2}{1}{Extended theory for certified finite-portfolio selection}{section.2}{PhysicaA_Supplement.pdf}}
\newlabel{S-sec:endpoint_selection}{{S2.1}{1}{Endpoint screening for a finite action portfolio}{subsection.2.1}{PhysicaA_Supplement.pdf}}
\newlabel{S-eq:portfolio_diameter}{{S1}{1}{Endpoint screening for a finite action portfolio}{equation.2.1}{PhysicaA_Supplement.pdf}}
\newlabel{S-eq:envelope_bet}{{S2}{1}{Endpoint screening for a finite action portfolio}{equation.2.2}{PhysicaA_Supplement.pdf}}
\newlabel{S-thm:endpoint_certificate}{{S1}{2}{Anytime endpoint certificate}{theorem.1}{PhysicaA_Supplement.pdf}}
\newlabel{S-eq:paired_bets}{{S3}{2}{A paired fallback with explicit unresolved outcomes}{equation.2.3}{PhysicaA_Supplement.pdf}}
\newlabel{S-thm:envelope_selection}{{S2}{2}{Envelope-first finite-portfolio selection}{theorem.2}{PhysicaA_Supplement.pdf}}
\newlabel{S-eq:portfolio_pac}{{S4}{2}{Envelope-first finite-portfolio selection}{equation.2.4}{PhysicaA_Supplement.pdf}}
\newlabel{S-alg:envelope_first}{{S1}{3}{A paired fallback with explicit unresolved outcomes}{algorithm.S1}{PhysicaA_Supplement.pdf}}
\newlabel{S-prop:zero_width_cost}{{S3}{3}{Zero-width endpoint cost}{theorem.3}{PhysicaA_Supplement.pdf}}
\newlabel{S-eq:zero_width_cost}{{S5}{3}{Zero-width endpoint cost}{equation.2.5}{PhysicaA_Supplement.pdf}}
\newlabel{S-sec:screen_cost}{{S2.3}{4}{When can endpoint screening repay its cost?}{subsection.2.3}{PhysicaA_Supplement.pdf}}
\newlabel{S-eq:chord_drift}{{S6}{4}{When can endpoint screening repay its cost?}{equation.2.6}{PhysicaA_Supplement.pdf}}
\newlabel{S-thm:screen_power}{{S5}{4}{A sufficient finite screening budget}{theorem.5}{PhysicaA_Supplement.pdf}}
\newlabel{S-eq:screen_power}{{S7}{4}{A sufficient finite screening budget}{equation.2.7}{PhysicaA_Supplement.pdf}}
\newlabel{S-prop:screen_drift}{{S6}{5}{Positive-drift screening cost}{theorem.6}{PhysicaA_Supplement.pdf}}
\newlabel{S-eq:screen_drift_cost}{{S8}{5}{Positive-drift screening cost}{equation.2.8}{PhysicaA_Supplement.pdf}}
\newlabel{S-prop:screen_break_even}{{S7}{5}{A sufficient copy-cost break-even condition}{theorem.7}{PhysicaA_Supplement.pdf}}
\newlabel{S-prop:screen_lower_bound}{{S8}{5}{An endpoint-only zero-width lower bound}{theorem.8}{PhysicaA_Supplement.pdf}}
\newlabel{S-eq:screen_lower_bound}{{S9}{5}{An endpoint-only zero-width lower bound}{equation.2.9}{PhysicaA_Supplement.pdf}}
\newlabel{S-eq:signed_screen_lower_bound}{{S10}{6}{An endpoint-only zero-width lower bound}{equation.2.10}{PhysicaA_Supplement.pdf}}
\newlabel{S-cor:implicit_region}{{S9}{6}{Implicit action regions and adaptive proposals}{theorem.9}{PhysicaA_Supplement.pdf}}
\newlabel{S-prop:race}{{S10}{7}{Stagewise approximate greedy selection}{theorem.10}{PhysicaA_Supplement.pdf}}
\newlabel{S-sec:witness_bound}{{S2.5}{7}{A fixed-witness bound for final-action certification}{subsection.2.5}{PhysicaA_Supplement.pdf}}
\newlabel{S-thm:witness}{{S11}{7}{Fixed-witness certification}{theorem.11}{PhysicaA_Supplement.pdf}}
\newlabel{S-sec:conditional_routing}{{S2.6}{8}{Independent-pilot routing and conditional risk reservation}{subsection.2.6}{PhysicaA_Supplement.pdf}}
\newlabel{S-thm:conditional_routing}{{S13}{8}{Conditional routing validity}{theorem.13}{PhysicaA_Supplement.pdf}}
\newlabel{S-sec:scoreguarantee}{{S3.1}{9}{A coefficient condition for the frozen score}{subsection.3.1}{PhysicaA_Supplement.pdf}}
\newlabel{S-eq:quadratic}{{S12}{9}{A coefficient condition for the frozen score}{equation.3.12}{PhysicaA_Supplement.pdf}}
\newlabel{S-eq:coefficients}{{S13}{9}{A coefficient condition for the frozen score}{equation.3.13}{PhysicaA_Supplement.pdf}}
\newlabel{S-thm:proxy}{{S14}{9}{Linear-exposure dominance}{theorem.14}{PhysicaA_Supplement.pdf}}
\newlabel{S-eq:dominance}{{S14}{9}{Linear-exposure dominance}{equation.3.14}{PhysicaA_Supplement.pdf}}
\newlabel{S-eq:kappabound}{{S15}{9}{Linear-exposure dominance}{equation.3.15}{PhysicaA_Supplement.pdf}}
\newlabel{S-cor:proxygreedy}{{S15}{9}{Exact greedy for the score}{theorem.15}{PhysicaA_Supplement.pdf}}
\newlabel{S-eq:proxyapprox}{{S16}{9}{Exact greedy for the score}{equation.3.16}{PhysicaA_Supplement.pdf}}
\newlabel{S-prop:localization}{{S16}{10}{Localization of positive score complementarity}{theorem.16}{PhysicaA_Supplement.pdf}}
\newlabel{S-sec:certificate}{{S3.2}{10}{Independent paired certification}{subsection.3.2}{PhysicaA_Supplement.pdf}}
\newlabel{S-eq:paired}{{S17}{10}{Independent paired certification}{equation.3.17}{PhysicaA_Supplement.pdf}}
\newlabel{S-thm:certificate}{{S17}{10}{One-query switching certificate}{theorem.17}{PhysicaA_Supplement.pdf}}
\newlabel{S-eq:eb}{{S18}{10}{One-query switching certificate}{equation.3.18}{PhysicaA_Supplement.pdf}}
\newlabel{S-eq:false}{{S19}{10}{One-query switching certificate}{equation.3.19}{PhysicaA_Supplement.pdf}}
\newlabel{S-eq:risk}{{S20}{10}{One-query switching certificate}{equation.3.20}{PhysicaA_Supplement.pdf}}
\newlabel{S-cor:portfolio}{{S18}{11}{Portfolio regret and time-uniform validity}{theorem.18}{PhysicaA_Supplement.pdf}}
\newlabel{S-eq:spending}{{S21}{11}{Portfolio regret and time-uniform validity}{equation.3.21}{PhysicaA_Supplement.pdf}}
\newlabel{S-tab:cohorts}{{S1}{11}{Evaluation cohorts, independently sampled units, and charged budgets. Final evaluation is separate from online selection}{table.1}{PhysicaA_Supplement.pdf}}
\newlabel{S-sec:conditional_details}{{S6.1}{12}{Calibration and independent confirmation}{subsection.6.1}{PhysicaA_Supplement.pdf}}
\newlabel{S-tab:robustness}{{S2}{13}{Post-hoc descriptive FEX robustness on the 60 confirmation graphs, averaging only the five eligible cells. No new inferential interval is attached}{table.2}{PhysicaA_Supplement.pdf}}
\newlabel{S-tab:grid}{{S3}{13}{Complete confirmation grid. Eligibility is fixed by independent calibration. Gain is NEX minus FEX; intervals adjust the 24 cellwise comparisons}{table.3}{PhysicaA_Supplement.pdf}}
\newlabel{S-sec:adaptive_study}{{S7}{14}{Independent-pilot routing and larger portfolios}{section.7}{PhysicaA_Supplement.pdf}}
\newlabel{S-tab:adaptive_cells}{{S4}{14}{Independent-pilot study with the frozen strong plain comparator. Every context is retained and the same exact extinction presolve is available to both methods. The common cap is 32,768 marginal copies}{table.4}{PhysicaA_Supplement.pdf}}
\newlabel{S-tab:adaptive_families}{{S5}{14}{All family-level plain/routed-screen wall ratios, reported descriptively without family-specific superiority tests}{table.5}{PhysicaA_Supplement.pdf}}
\newlabel{S-fig:adaptive_comparison}{{S1}{15}{Independent-pilot routing. (a) Mean charged copies and certified-query counts, including the common exact presolve. (b) The pooled planned plain/routed-screen geometric wall-time ratio with its conservative 97.5\% interval, and family-level descriptive points without intervals. Values above one favor routed screening. All pilot costs are included}{figure.1}{PhysicaA_Supplement.pdf}}
\newlabel{S-fig:pool_sensitivity}{{S2}{17}{Candidate-domain sensitivity, separated from computational scaling. Points show NEX-minus-FEX loss for all six size/pool settings, with the original eight-endpoint-family adjusted $99.375\%$ intervals. Distinct pool definitions are categorical and are not joined by lines. All-node pools contain $n$ candidates. Only the $n=300$, all-node interval excludes zero}{figure.2}{PhysicaA_Supplement.pdf}}
\newlabel{S-fig:complete_search_mechanisms}{{S5}{22}{Prespecified diagnostic workloads. (a) Native time on identical recorded query portfolios. (b) All candidate-count and overlap workloads; the two portfolio constructions also differ in their constituent sets, so these are not isolated causal effects of overlap}{figure.5}{PhysicaA_Supplement.pdf}}
\newlabel{S-fig:portfolio}{{S6}{23}{Portfolio evaluation at unchanged graphical joint law. (a) Graph-level geometric mean CPU-time ratios with three-comparison adjusted 98.333\% bootstrap intervals; ratios above one favor local adaptive batching. (b) Every family/size/portfolio replay-to-local ratio, shown descriptively without per-cell significance claims. Incremental RTC is a distinct coupling; its runtime comparison is at equal marginal-continuation budget}{figure.6}{PhysicaA_Supplement.pdf}}
\newlabel{S-sec:envelope_experiment}{{S18.2}{23}{Order-informed portfolio selection}{subsection.18.2}{PhysicaA_Supplement.pdf}}
\newlabel{S-fig:batch_scaling}{{S7}{24}{Descriptive portfolio scaling on fresh larger graphs. Lines show median CPU seconds per marginal continuation over eight graph instances per size; shading denotes their interquartile range, not a confidence interval. All methods evaluate the same sets and two tapes per workload. These randomized augmented states test computation and are not naturally generated histories}{figure.7}{PhysicaA_Supplement.pdf}}
\newlabel{S-tab:envelope_selection}{{S14}{25}{Fresh finite-portfolio selection: certification, endpoint-only returns, total charged online copies and wall time. Means are over all 96 contexts, including unresolved cases}{table.14}{PhysicaA_Supplement.pdf}}
\newlabel{S-fig:envelope_selection}{{S8}{25}{Order-informed selection on the fresh locked cohort. (a) Comparator/envelope-first wall-time ratios; pooled planned contrasts use adjusted 97.5\% intervals, and family intervals are descriptive 95\% intervals. Values above one favor envelope-first selection. (b) Every context's total charged online copies; auxiliary endpoints count in full. Costs include unresolved runs; neither planned comparison meets the matched-certified-quality gate. These are fixed-portfolio tolerance certificates, not guarantees of globally optimal intervention}{figure.8}{PhysicaA_Supplement.pdf}}
\newlabel{S-sec:factorial}{{S18.3}{25}{Separating endpoint information from coupling implementation}{subsection.18.3}{PhysicaA_Supplement.pdf}}
\newlabel{S-tab:factorial_cells}{{S15}{26}{Independent engine/screen study. Means include unresolved queries; endpoint-only certificates do not simulate the individual candidates}{table.15}{PhysicaA_Supplement.pdf}}
\newlabel{S-tab:factorial_cost}{{S16}{26}{All four planned graph-level geometric capped-time contrasts on the new cohort. Ratios above one favor the denominator; intervals are adjusted 98.75\% intervals. Graph signs refer to mean log ratios. Certification gates must also pass for an equal-quality efficiency claim}{table.16}{PhysicaA_Supplement.pdf}}
\newlabel{S-fig:factorial_selection}{{S9}{26}{Separating screening from the simulation engine. (a) All 96 query costs per cell; horizontal bars are medians, red crosses denote unresolved outputs, and labels give certification counts. (b) All four planned geometric wall-time ratios and adjusted 98.75\% intervals. Costs include unresolved outputs and cannot alone establish equal-quality acceleration}{figure.9}{PhysicaA_Supplement.pdf}}
\newlabel{S-sec:witness_study}{{S18.4}{27}{Independent evaluation of the fixed-witness refinement}{subsection.18.4}{PhysicaA_Supplement.pdf}}
\newlabel{S-tab:witness_cells}{{S17}{27}{Fixed-witness refinement on the independently planned third cohort. Means include every query; the common cap is 32,768 marginal copies}{table.17}{PhysicaA_Supplement.pdf}}
\newlabel{S-fig:witness_comparison}{{S10}{28}{Fixed-witness refinement. (a) Mean charged copies and certified-query counts. (b) Both planned geometric wall-time ratios with adjusted 97.5\% intervals. Values above one favor the denominator. Certification counts and costs are reported separately; the common cap exceeds the earlier studies' cap}{figure.10}{PhysicaA_Supplement.pdf}}
\newlabel{S-fig:complete_search}{{S11}{30}{Complete adaptive search. (a) The primary replay/local ratio with its $95\%$ interval and two secondary ratios with adjusted $97.5\%$ intervals. The open RTC marker denotes equal budget, rather than matched selection quality. (b) All graph-level replay/local ratios; blue and coral indicate sizes $100$ and $300$, and horizontal bars show descriptive group geometric means}{figure.11}{PhysicaA_Supplement.pdf}}
\newlabel{S-fig:optimized_replay}{{S12}{30}{Persistent-topology replay confirmation. (a) Both planned whole-search ratios with Bonferroni-adjusted $97.5\%$ intervals. (b) All 32 independent graph-level persistent-replay/local ratios, grouped by family and node count. The dashed reference denotes equal time; plotted group means are descriptive}{figure.12}{PhysicaA_Supplement.pdf}}
\newlabel{S-sec:design_details}{{S24.1}{30}{Complete history and parameter specification}{subsection.24.1}{PhysicaA_Supplement.pdf}}
\newlabel{S-sec:search_details}{{S24.2}{31}{Frozen proposal score and search specification}{subsection.24.2}{PhysicaA_Supplement.pdf}}
\newlabel{S-eq:score}{{S22}{31}{Frozen proposal score and search specification}{equation.24.22}{PhysicaA_Supplement.pdf}}
\newlabel{S-sec:portfolio_details}{{S24.3}{32}{Order-informed certificates and their role}{subsection.24.3}{PhysicaA_Supplement.pdf}}
\newlabel{S-eq:portfolio_details}{{S23}{32}{Order-informed certificates and their role}{equation.24.23}{PhysicaA_Supplement.pdf}}
\newlabel{S-sec:variance_details}{{S24.4}{32}{Fixed-contrast RTC comparison}{subsection.24.4}{PhysicaA_Supplement.pdf}}
\newlabel{S-sec:scaling_details}{{S24.5}{32}{Full scaling and pool sensitivity}{subsection.24.5}{PhysicaA_Supplement.pdf}}
\newlabel{S-fig:scalability}{{S13}{33}{Computational scaling. (a) Per-trajectory wall-time medians and interquartile ranges over sixteen measured graph instances per size, for graphical coupling and incremental RTC. Shading is an interquartile range, not a confidence interval. (b) Graphical wall time versus the measured proposal-work proxy, showing every scaling observation; the dashed descriptive log--log fit has slope $1.003$. The figure characterizes measured computational work rather than a uniform runtime advantage. Candidate-pool search effects are reported separately in Supplementary Fig.~\ref {fig:pool_sensitivity}}{figure.13}{PhysicaA_Supplement.pdf}}
\newlabel{S-sec:transfer_details}{{S24.6}{33}{Strong baselines and observed-group transfer}{subsection.24.6}{PhysicaA_Supplement.pdf}}
\newlabel{S-sec:certification_details}{{S24.7}{33}{Optional certification and charged cost}{subsection.24.7}{PhysicaA_Supplement.pdf}}
\newlabel{S-sec:quality_estimates}{{S24.8}{33}{Additional fixed-quality loss estimates}{subsection.24.8}{PhysicaA_Supplement.pdf}}
\newlabel{S-sec:memory_response_details}{{S25}{33}{Controlled-initial-memory response experiment}{section.25}{PhysicaA_Supplement.pdf}}
\newlabel{S-fig:original_memory_contexts}{{S14}{35}{Separate contexts underlying main Fig.~\ref {fig:memory_response}. Each row is one fixed graph and each column a face-strength setting. Teal and red denote initial $-T_0$ and $+T_0$; solid and dashed lines denote no control and four targets. Curves average 128 tapes within each context. The vertical line marks withdrawal. Node trajectories for the two memory initializations coincide at $s=0$}{figure.14}{PhysicaA_Supplement.pdf}}
\newlabel{S-sec:response_extension}{{S26}{35}{Memory-mechanism and comparison-bound diagnostics}{section.26}{PhysicaA_Supplement.pdf}}
\newlabel{S-tab:memory_modes}{{S20}{36}{Mean occupation loss across the eight fixed contexts. Dynamic-setting labels give initial memory and $\lambda $. All five controls are included}{table.20}{PhysicaA_Supplement.pdf}}
\newlabel{S-fig:bound_diagnostics}{{S15}{37}{All 72 empirical variance-to-bound ratios for nested and nonnested comparisons. Colors and symbols identify families. The horizontal dotted line is equality with the plug-in bound. Ratios describe the empirical distribution of tapes; they do not certify population tightness}{figure.15}{PhysicaA_Supplement.pdf}}
\newlabel{S-fig:memory_contexts}{{S16}{37}{Initial-memory loss gaps for every graph. Red: no control; teal: four targets. Solid and dashed lines distinguish the two graphs in each family. Every estimated graph-specific gap decreases over the three tested rates; no continuous-parameter monotonicity is inferred}{figure.16}{PhysicaA_Supplement.pdf}}
\newlabel{S-sec:threshold_sensitivity}{{S27}{37}{Post-hoc sensitivity to the selection-quality threshold}{section.27}{PhysicaA_Supplement.pdf}}
\newlabel{S-tab:threshold_sensitivity}{{S21}{38}{Post-hoc stable-crossing budgets in physical trajectories. Pooled means use 32 graphs; family means use eight. Thresholds and family summaries here carry no new inferential claims}{table.21}{PhysicaA_Supplement.pdf}}
\newlabel{S-sec:runtime_implementation}{{S28.1}{38}{Runtime accounting and persistent replay}{subsection.28.1}{PhysicaA_Supplement.pdf}}
\newlabel{S-sec:search_implementation}{{S28.2}{39}{Candidate screening and search}{subsection.28.2}{PhysicaA_Supplement.pdf}}
\newlabel{S-sec:natural_memory_details}{{S29}{39}{Natural-history memory and feedback study}{section.29}{PhysicaA_Supplement.pdf}}
\newlabel{S-tab:natural_memory_all}{{S22}{39}{All natural-history loss contrasts. Positive values mean the first rule has greater occupation loss. F: dynamic natural minus frozen; Z: dynamic natural minus zero; R: dynamic natural minus reset dynamic. Control 0 is absent and 1 uses four targets. Signs count graph-level means}{table.22}{PhysicaA_Supplement.pdf}}
\newlabel{S-sec:full_computational_application}{{S30}{40}{Simulation implementation and search applications}{section.30}{PhysicaA_Supplement.pdf}}
\newlabel{S-prop:work}{{S19}{40}{Work bound and exact extinction stopping}{theorem.19}{PhysicaA_Supplement.pdf}}
\newlabel{S-eq:work}{{S24}{40}{Work bound and exact extinction stopping}{equation.30.24}{PhysicaA_Supplement.pdf}}
\newlabel{S-alg:simulation}{{S2}{41}{Simulation implementation and search applications}{algorithm.S2}{PhysicaA_Supplement.pdf}}
\newlabel{S-prop:unbiased}{{S20}{41}{Fixed-query unbiasedness}{theorem.20}{PhysicaA_Supplement.pdf}}
\newlabel{S-alg:batch}{{S3}{42}{Batch evaluation of candidate actions}{algorithm.S3}{PhysicaA_Supplement.pdf}}
\newlabel{S-sec:search}{{S30.2}{42}{Search application}{subsection.30.2}{PhysicaA_Supplement.pdf}}
\newlabel{S-eq:portfolio_main}{{S25}{42}{Order-informed certification}{equation.30.25}{PhysicaA_Supplement.pdf}}
\newlabel{S-alg:search}{{S4}{43}{Anchor, score, and exchange neighborhoods}{algorithm.S4}{PhysicaA_Supplement.pdf}}
\newlabel{S-fig:cost}{{S17}{43}{Cost and selection quality for four equal-size nonnested actions. Panels (a) and (b) compare physical trajectory budgets and wall time on a shared vertical scale. Insets enlarge the original high-budget Graphical/RTC and thinning/antithetic comparisons. Markers alternate across observed budgets to reduce occlusion; all mean values and interval endpoints are reported. The reference is chosen on independent sample A and assessed on sample B. Trajectory counts charge both antithetic partners. The fixed threshold is excess loss $0.02$; curve error bars are descriptive 95\% intervals, while cost ratios use the adjusted intervals in Table~\ref {tab:cost}. Wall time covers online selection after candidates are fixed. The zero-simulation R-degree anchor is reported numerically below the panels and lies above their common vertical range}{figure.17}{PhysicaA_Supplement.pdf}}
\newlabel{S-tab:cost}{{S23}{44}{Cost at the fixed quality target: comparator cost divided by graphical cost. All intervals are Bonferroni-adjusted 99.375\% intervals for the eight-endpoint family}{table.23}{PhysicaA_Supplement.pdf}}
\newlabel{S-fig:jcs_runtime}{{S18}{45}{Complete-search computation and graph-level consistency. (a) The 64-graph study reports its primary replay/local ratio with a $95\%$ interval and secondary ratios with adjusted $97.5\%$ intervals. The open RTC marker denotes equal marginal budget and a different coupling, without matched-quality evidence. (b) The 32-graph confirmation uses persistent topology and one full-prefix call per action and stage; both intervals are adjusted $97.5\%$. (c,d) Every graph-level primary ratio is shown, distinguished by network size; short vertical bars are descriptive family geometric means. Values above one favor local batching. Primary ratios correspond to geometric-ratio time savings of 75.7\% and 73.6\%. Each cohort uses its own multiple-comparison adjustment}{figure.18}{PhysicaA_Supplement.pdf}}
\newlabel{S-fig:confirmation}{{S19}{46}{Where the conditional search gain occurs. All methods use $1,024$ online continuations per policy, with independent final evaluation. (a,b) The complete eight-cell confirmation grid on 60 graphs: entries are baseline loss minus FEX loss, averaged over histories and then within each family or across all graphs. Both panels use the same zero-centered linear color scale. Filled markers identify the five cells selected by independent calibration; open markers identify the other cells. Here $s$ is the face-strength scale and $u=3$ throughout. Family cells and pooled cell means are descriptive; the separately adjusted cellwise intervals remain in Supplementary Table~\ref {tab:grid}. (c) Aggregation over the five calibrated cells. Aggregate primary FEX--NEX and secondary FEX--greedy intervals are adjusted $97.5\%$ intervals for the two planned comparisons; family intervals are descriptive $95\%$ intervals. Histories and cells stay grouped within each graph. Positive differences favor FEX}{figure.19}{PhysicaA_Supplement.pdf}}
\newlabel{S-fig:gridheatmap}{{S19}{46}{Where the conditional search gain occurs. All methods use $1,024$ online continuations per policy, with independent final evaluation. (a,b) The complete eight-cell confirmation grid on 60 graphs: entries are baseline loss minus FEX loss, averaged over histories and then within each family or across all graphs. Both panels use the same zero-centered linear color scale. Filled markers identify the five cells selected by independent calibration; open markers identify the other cells. Here $s$ is the face-strength scale and $u=3$ throughout. Family cells and pooled cell means are descriptive; the separately adjusted cellwise intervals remain in Supplementary Table~\ref {tab:grid}. (c) Aggregation over the five calibrated cells. Aggregate primary FEX--NEX and secondary FEX--greedy intervals are adjusted $97.5\%$ intervals for the two planned comparisons; family intervals are descriptive $95\%$ intervals. Histories and cells stay grouped within each graph. Positive differences favor FEX}{figure.19}{PhysicaA_Supplement.pdf}}
\newlabel{S-tab:jcs_cohorts}{{S24}{47}{Principal evaluation cohorts. A continuation is one physical trajectory. Inference groups all histories and settings of each graph. Supplementary Table~\ref {tab:cohorts} lists all additional cohorts and budgets}{table.24}{PhysicaA_Supplement.pdf}}
\newlabel{S-sec:earlier_memory_cohorts}{{S31}{47}{Earlier controlled-memory cohorts and full design}{section.31}{PhysicaA_Supplement.pdf}}
\newlabel{S-sec:verification}{{S31.2}{48}{Exactness and mathematical verification}{subsection.31.2}{PhysicaA_Supplement.pdf}}
\newlabel{S-sec:memory_response}{{S31.3}{48}{Controlled initial memory and spreading response}{subsection.31.3}{PhysicaA_Supplement.pdf}}
\newlabel{S-sec:timescales}{{S31.4}{48}{Static bias, dynamic memory and time scales}{subsection.31.4}{PhysicaA_Supplement.pdf}}
\newlabel{S-fig:memory_response}{{S20}{49}{Controlled-initial-memory response (descriptive). Columns show $s=0$ and $s=1$; rows show mean rumor occupation and mean face memory. Curves average $128$ trajectories on each of four fixed graph/initial-state contexts. Color denotes initial face memory; line style denotes intervention. Shading shows conditional Monte Carlo standard errors, not uncertainty over a graph population. The vertical line marks control withdrawal at $D=2$. The two memory-initialization node curves coincide when $s=0$. All conditions are shown}{figure.20}{PhysicaA_Supplement.pdf}}
\newlabel{S-sec:natural_memory}{{S31.5}{49}{Naturally generated memories and feedback}{subsection.31.5}{PhysicaA_Supplement.pdf}}
\newlabel{S-fig:timescales}{{S21}{50}{Memory mechanisms at fixed higher-order transmission (descriptive; eight graphs). Red denotes no control and teal denotes four targets. (a) Loss difference between initial $+T_0$ and $-T_0$ under dynamic memory; dotted lines show frozen-memory differences. Error bars are one conditional Monte Carlo standard error for the equally weighted graph average; they quantify tape noise with these graphs fixed, not uncertainty across a population of graphs. (b) Mean loss from initial $-T_0$ (solid circles) and $+T_0$ (dashed squares); dotted horizontal lines show memory fixed at zero. Both signs use identical initial node states. Frozen memory means $\lambda =0$; the other tick labels give the relaxation rate. Lines connect tested settings, with no interpolation-based inference}{figure.21}{PhysicaA_Supplement.pdf}}
\newlabel{S-fig:natural_memory}{{S22}{51}{Naturally generated memory (32 graphs, 64 histories; descriptive). (a) Post-control occupation under four memory rules, with four targets. Frozen and zero controls are identical across the continuation-rate grid. (b) Spatial mean-memory contrast between dynamic natural and dynamic reset branches after subtracting direct exponential decay of the initial memory; the vertical line marks control withdrawal. Full pointwise uncertainty and both control conditions are reported in the supplement and accompanying CSV files}{figure.22}{PhysicaA_Supplement.pdf}}
\newlabel{S-sec:physics_new}{{S32}{51}{Withdrawal-time experiments and observed contact data}{section.32}{PhysicaA_Supplement.pdf}}
\newlabel{S-tab:withdrawal_all}{{S25}{52}{Complete withdrawal-time results. $G$ is total loss reduction, $G_A$ its memory component, and $\eta _A$ the memory fraction. The interval applies to $\eta _A$ in percent}{table.25}{PhysicaA_Supplement.pdf}}
\newlabel{S-sec:triangle_new}{{S33}{53}{Single-triangle retention and rebound}{section.33}{PhysicaA_Supplement.pdf}}
\newlabel{S-tab:triangle_all}{{S26}{53}{All triangle cells. $\mathcal R$ is mean sampled pathwise rebound; $\mathcal R_{\rm mean}$ is the positive overshoot of the mean curve; memory fractions are in percent}{table.26}{PhysicaA_Supplement.pdf}}
\newlabel{S-fig:triangle_ensemble}{{S23}{54}{Overshoot of the mean occupation curve on the same complete triangle grid. Its scale is much smaller than mean within-path rebound}{figure.23}{PhysicaA_Supplement.pdf}}
\newlabel{S-sec:hyperedge4}{{S34}{54}{Weighted four-uniform factors}{section.34}{PhysicaA_Supplement.pdf}}
\newlabel{S-fig:hyperedge4}{{S24}{55}{Ordered occupation response on weighted four-uniform factors under two bounded antitone drives. The vertical line marks withdrawal}{figure.24}{PhysicaA_Supplement.pdf}}
\newlabel{S-sec:precision_new}{{S35}{55}{Threshold-free precision and cost}{section.35}{PhysicaA_Supplement.pdf}}
}{\externaldocument[S-]{supplement}[PhysicaA_Supplement.pdf]}
\usepackage[hidelinks]{hyperref}
\providecommand{\urlprefix}{}
\journal{Physica A: Statistical Mechanics and its Applications}
\date{31 October 2026}
% Keep wide result figures near their first discussion and allow text below.
\renewcommand{\dbltopfraction}{0.90}
\renewcommand{\dblfloatpagefraction}{0.85}
\renewcommand{\textfraction}{0.08}
\setcounter{dbltopnumber}{2}
\theoremstyle{plain}
\newtheorem{theorem}{Theorem}
\newtheorem{proposition}[theorem]{Proposition}
\newtheorem{corollary}[theorem]{Corollary}
\newtheorem{lemma}[theorem]{Lemma}
\newcommand{\ind}{\mathbf{1}}
\newcommand{\E}{\mathbb{E}}
\newcommand{\Prob}{\mathbb{P}}
\newcommand{\Fset}{\mathcal{F}}
\newcommand{\Var}{\operatorname{Var}}
\newcommand{\Gain}{\mathcal{G}}
\hypersetup{pdftitle={Group memory and post-intervention response in higher-order spreading},pdfauthor={Junkai Fan}}
\begin{document}
\begin{frontmatter}
\title{Group memory and post-intervention response in higher-order spreading}
\author[neu]{Junkai Fan\corref{cor1}}
\ead{fanjk@mails.neu.edu.cn}
\cortext[cor1]{Corresponding author.}
\affiliation[neu]{organization={School of Business Administration, Northeastern University},city={Shenyang},postcode={110169},country={China}}
\begin{abstract}
With prescribed synthetic transition rates, group memory can preserve part of a temporary intervention's effect after the intervention ends. We study competing spreading on higher-order networks, where recurrent node transitions interact with continuously relaxing signed face memories. At withdrawal, a counterfactual comparison keeps the intervention's terminal node state and replaces its memory with that of a paired untreated process. This separates node-state and memory contributions to subsequent occupation loss. Simulations on synthetic networks and sensor-derived high-school and primary-school contact topologies show nonzero memory-mediated contributions under the tested settings. An exact order-preserving construction combines an ordering of the three node states with reverse numerical memory order and accommodates the deterministic memory flow. It yields pathwise monotonicity under ordered controls, nonnegative nested loss differences and intersection--union bounds for nonnested comparisons. Classical variance inequalities then supply conservative fluctuation and sufficient sampling bounds. A complete three-node example has strict pair complementarity and a zero submodularity ratio, separating monotonicity from diminishing returns. Small-system calculations distinguish a smooth direct-memory retention scale from sampled post-withdrawal rebound; naturally generated histories further separate forgetting of initial memory from state-driven feedback. These are finite-time results, not evidence for a phase boundary. A supplementary computational application preserves the joint trajectory law while reducing evaluation cost in the tested workloads.
\end{abstract}
\begin{keyword}
Higher-order networks \sep Competing spreading \sep Group memory \sep Monotone coupling \sep Stochastic simulation \sep Intervention response
\end{keyword}
\end{frontmatter}
\raggedbottom
\emergencystretch=1.5em
\section{Introduction}
Temporary correction can alter the visible state of a spreading system before it removes the influence stored in its groups. Once the external input ends, the same node configuration can respond differently according to that hidden state. Higher-order transmission makes this distinction especially relevant: reinforcement depends on several donors acting together, so the past activity of a group can affect the propagation channels available after withdrawal~\cite{iacopini2019,battiston2020,battiston2021}. We study this finite-time response in a recurrent interacting system with continuously relaxing face memory.

Three questions organize the work. First, prescribed positive and negative memories demonstrate sensitivity to a hidden state, but do not isolate the memory produced by an intervention itself. That contribution requires a withdrawal-time comparison holding the intervention's terminal node state fixed. Second, the direction of an intervention response and its diminishing returns are different properties. Classical influence guarantees exploit submodularity~\cite{kempe2003,nemhauser1978}, whereas group reinforcement permits complementary targets. For the present memory-dependent process, an exact comparison must handle stochastic state changes and continuous memory flow together. Third, evaluating several candidate perturbations requires preserving the dependence used to compare them; exact marginal simulations alone do not determine a joint comparison law.

Our model has rumor, inactive and correction states, with recurrent switching. A signed memory on each filled triangle is driven by its local majority and relaxes toward that drive. Temporary correction increases the adoption intensity at selected nodes. The observable is cumulative rumor occupation during a fixed interval after withdrawal. The continued-influence effect motivates distinguishing a present correction state from a persisting trace of earlier endorsement~\cite{lewandowsky2012}. Here that distinction is represented at group level through specified rates. Sensor-derived contact groups from a high school and a primary school provide two observed topologies for the same synthetic dynamics~\cite{benson2018,highschool2015,primaryschool2011}.

The contributions correspond to the three questions:
\begin{enumerate}
\item A withdrawal-time memory comparison separates the intervention benefit into a node-state component and a memory component. Controlled simulations on synthetic and observed contact topologies quantify this distinction and its dependence on intervention duration.
\item An augmented-state order combines $R\prec U\prec C$ with reverse numerical memory order. An exact common-random-event construction gives pathwise comparisons, fluctuation bounds and intersection--union envelopes. A complete-process counterexample shows that monotone response can coexist with strict complementarity and a zero submodularity ratio.
\item An exact joint evaluator shares computation while preserving the chosen trajectory law. Its computational consequences are assessed separately from the physical results, including threshold-free error and variance measurements and the limits of improvement over common reaction paths.
\end{enumerate}

The order is model-specific. An ordering of the node labels alone does not check whether the deterministic memory flow and reset intensities remain compatible. We verify that compatibility for the reverse signed-memory order, including sufficient conditions for weighted higher-order factors and heterogeneous relaxation. This is the step linking the continuous memory to attractive-process comparison; the use of common events is established methodology.

\subsection{Relation to previous work}
Higher-order contagion models show how nonlinear donor requirements change spreading relative to pairwise interaction~\cite{iacopini2019,neuhauser2020}. Reviews describe their structural and dynamical settings~\cite{battiston2020,battiston2021,bick2023,boccaletti2023}. Temporal contact and path-dependent network models represent additional forms of history dependence~\cite{holme2012,lambiotte2019}; here a continuous variable is attached to each fixed group and remains active between node transitions. We use the timestamped contact simplices assembled for simplicial-closure analysis~\cite{benson2018}. These groups are maximal cliques within sensor time windows, rather than triangles filled after aggregating all contacts.

Graphical constructions for attractive multitype processes compare ordered paths using common random maps~\cite{stover2010}. The present augmented state is a piecewise-deterministic Markov process~\cite{davis1984}. Exact reaction simulation and event thinning provide its marginal construction~\cite{gillespie1977,lewis1979}; common operational-time reaction paths offer another coupling~\cite{anderson2007,anderson2012}. Anderson and Yuan develop stacked coupling for time-dependent reaction intensities~\cite{andersonyuan2019}. Anderson and Ma's recent author preprint extends stacked coupling to finite collections of reaction-network paths~\cite{andersonma2026}. Those constructions establish shared-event simulation. Our contribution is to verify a joint state--memory order for the signed-memory flow and derive its consequences for intervention response. Attractiveness additionally requires compatibility of the node updates and continuous memory flow.

Nonlinear threshold diffusion and higher-order seeding need not have diminishing returns~\cite{schoenebeck2022,zhang2024,qu2025}. Weak-submodularity guarantees explicitly require a positive ratio~\cite{das2018}. Our three-node example tests that requirement in the complete recurrent process, including forgetting, temporary control and memory relaxation. Search heuristics, influence-blocking comparisons and selection certificates are treated as supplementary applications rather than as a universal optimization claim.

\section{Memory-dependent competing spreading}
\subsection{Network, states, and signed memory}
Let $G=(V,E,\Fset)$ be a finite simplicial network with $n=|V|\ge1$ nodes, undirected pairwise edges $E$, and filled triangles $\Fset$. Every filled triangle contains its three constituent edges. A three-cycle in the graph is not automatically filled. At time $t$, node $i$ has state $X_i(t)\in\{R,U,C\}$, representing rumor transmission, inactivity, or correction transmission. Write $x=(x_i)_{i\in V}$ for a node configuration.

Each face $\sigma\in\Fset$ has signed memory $A_\sigma(t)\in[-T_0,T_0]$, where $0<T_0\le1$. Positive memory favors rumor reinforcement and negative memory favors correction reinforcement. This is a minimal group-level representation of persisting influence, motivated by the continued-influence effect~\cite{lewandowsky2012}; its transition rates are prescribed synthetic parameters. With $N_R(\sigma,x)$ and $N_C(\sigma,x)$ the corresponding state counts, define
\begin{align}
 \Phi_\sigma(x)&=T_0\big[\ind\{N_R(\sigma,x)\ge2\}
                      -\ind\{N_C(\sigma,x)\ge2\}\big],\label{eq:drive}\\
 \dot A_\sigma(t)&=\lambda[\Phi_\sigma(X(t))-A_\sigma(t)],
 \qquad \lambda\ge0.\label{eq:memory}
\end{align}
Time is measured in units of the inactive-reset rate, which is one. The memory $A_\sigma$ and its scale $T_0$ are dimensionless; $\beta_1$, $\beta_2$, $u$ and $\lambda$ are rates in these units. The parameter $T_0$ sets both the transmission prefactor and memory amplitude. Consequently, varying $T_0$ alone changes both effects. The numerical memory comparisons below hold $T_0$ fixed, with $\beta_1T_0$ and $\beta_2T_0$ normalized by mean edge degree and face participation. Since $|\Phi_\sigma|\le T_0$, the interval $[-T_0,T_0]$ is invariant. While $X$ remains constant,
\begin{equation}
 A_\sigma(t+h)=\Phi_\sigma(X(t))+
 [A_\sigma(t)-\Phi_\sigma(X(t))]e^{-\lambda h}.\label{eq:flow}
\end{equation}

For nonnegative pairwise and face rates $\beta_1,\beta_2$, the adoption pressures are
\begin{align}
 \ell_i^R(x,a)&=\beta_1T_0\sum_{j:\{i,j\}\in E}\ind\{x_j=R\}\notag\\
 &\quad+\beta_2T_0\sum_{\sigma\ni i}(1+a_\sigma)
       \prod_{j\in\sigma\setminus\{i\}}\ind\{x_j=R\},\label{eq:rumor}\\
 \ell_i^C(x,a,u)&=\beta_1T_0\sum_{j:\{i,j\}\in E}\ind\{x_j=C\}\notag\\
 &\quad+\beta_2T_0\sum_{\sigma\ni i}(1-a_\sigma)
       \prod_{j\in\sigma\setminus\{i\}}\ind\{x_j=C\}+u_i(t).
 \label{eq:correction}
\end{align}
All pressures are nonnegative. Table~\ref{tab:transitions} gives the complete transition law. In particular, active nodes can switch directly between $R$ and $C$; both active states return to $U$ at unit rate. The full state $Z=(X,A)$ is a piecewise-deterministic Markov process~\cite{davis1984}.

\begin{table}[t]
\centering
\caption{Complete node-transition law at the current augmented state.}
\label{tab:transitions}
\begin{tabular}{ll}
\toprule
Transition & Intensity\\
\midrule
$U\to R$, $C\to R$ & $\ell_i^R(x,a)$\\
$U\to C$, $R\to C$ & $\ell_i^C(x,a,u)$\\
$R\to U$, $C\to U$ & $1$\\
\bottomrule
\end{tabular}
\end{table}

Write $x^{i:v}$ for $x$ with coordinate $i$ reset to $v$. On functions differentiable in memory, the augmented generator is
\begin{align}
 \mathcal A_t f(x,a)
 &=\sum_{\sigma\in\Fset}\lambda[\Phi_\sigma(x)-a_\sigma]
              \partial_{a_\sigma}f(x,a)\notag\\
 &\quad+\sum_{i\in V}\sum_{v\in\{R,C,U\}}q_i^v(x,a,t)
      [f(x^{i:v},a)-f(x,a)],\label{eq:generator}
\end{align}
where $q_i^R=\ell_i^R$, $q_i^C=\ell_i^C$, and $q_i^U=1$. The reset decomposition automatically makes self-transitions inert. It is also why forgetting is compatible with the three-state order: one common reset map represents both active-to-inactive transitions.

Current node labels need not be a Markov state. For example, two augmented states with the same face labels $(R,R,U)$ but different memories give different instantaneous $U\to R$ rates when $\beta_2>0$. Simulation therefore includes memory in the state.

\subsection{Action, loss, and complementarity}
An intervention selects $S\subseteq V$ with $|S|=k$, $1\le k\le n$, and applies intensity $u\ge0$ per selected node for a fixed duration $D\ge0$:
\begin{equation}
 u_i^S(t)=u\ind\{i\in S\}\ind\{0\le t<D\}.\label{eq:control}
\end{equation}
Each feasible set has intensity--time cost $kuD$. The start time is fixed; $D$ is a duration. For a post-intervention horizon $H>0$ and a fixed initial context $(G,x,a)$, define
\begin{align}
 L(S,\xi)&=\int_D^{D+H}\frac1n\sum_{i\in V}
                 \ind\{X_i^S(t;\xi)=R\}\,dt,\label{eq:loss}\\
 Q(S)&=\E[L(S,\xi)\mid G,x,a].\label{eq:objective}
\end{align}
Here $\xi$ denotes one random simulation tape and $0\le L\le H$. The instantaneous occupation is $\rho_R(t)=n^{-1}\sum_i\ind\{X_i(t)=R\}$; its integral $L$ has time units, while $L/H$ is the average occupation fraction over the observation window. The benefit of a temporary perturbation is $g(S)=Q(\varnothing)-Q(S)$. Fixed-cardinality search for a small-loss target set is a secondary application reported in the supplement. For $i,j\notin B_g$, $i\ne j$, let
\begin{align}
 \Delta_i(B_g)&=g(B_g\cup\{i\})-g(B_g),\\
 \Gamma_{ij}(B_g)&=g(B_g\cup\{i,j\})-g(B_g\cup\{i\})\notag\\
 &\quad-g(B_g\cup\{j\})+g(B_g).\label{eq:complementarity}
\end{align}
Positive $\Gamma_{ij}$ means that the two interventions are complementary at background $B_g$. This diagnostic compares perturbations of different sizes; fixed-cardinality search comparisons are reported separately.

Figure~\ref{fig:overview} summarizes how memory enters the dynamics and how common reset events compare alternative controls.
\begin{figure*}[t]
\centering\includegraphics[width=.98\textwidth]{mechanism_reference_vector.pdf}
\caption{Augmented memory, ordered coupling and joint evaluation. (a) Identical node states can carry different signed memories. The displayed $R/U/C$ face has $\Phi_\sigma=0$, so its memories relax toward zero. (b) Common destination-reset proposals preserve $R\prec U\prec C$ and reverse numerical memory order under nested controls. Bars show schematic acceptance regions; the common $U$ reset is omitted. Each proposal uses a fresh mark, shared across actions. (c) For $S_1=\{1,2,3\}$ and $S_2=\{3,4,5\}$, intersection $I$ and union $J$ bound the loss difference without ordering the actions. The auxiliary union may exceed $k$. (d) Replay and batching preserve the full joint law; fresh tapes extend invariance to adaptive search. The diagram is schematic.}
\label{fig:overview}
\end{figure*}

\section{Exact ordered coupling and its consequences}
\label{sec:theory}
\subsection{An order compatible with signed memory}
Order node states by $R\prec U\prec C$. For augmented states $z=(x,a)$ and $z'=(x',a')$, define
\begin{equation}
 z\preceq z'\quad\Longleftrightarrow\quad
 x_i\preceq x_i'\ (i\in V),\qquad
 a_\sigma\ge a_\sigma'\ (\sigma\in\Fset).
 \label{eq:order}
\end{equation}
The opposite direction on memory is essential: an increase toward correction in the node coordinates drives signed endorsement toward negative values. Equation~\ref{eq:order} defines the order used throughout the analysis.

\begin{theorem}[Exact grand coupling and control comparison]
\label{thm:coupling}
Fix a finite topology, $0<T_0\le1$, $\lambda\ge0$, nonnegative finite transmission rates, and a finite horizon $K_{\mathrm{hor}}$. Let a finite collection of deterministic measurable controls be bounded by common per-node envelopes. There is a common marked Poisson construction with the correct marginal process for every control. Whenever two initial states and controls satisfy
\begin{equation}
 z_0\preceq z'_0,\qquad u_i(t)\le u'_i(t)
 \quad(i\in V,\ 0\le t\le K_{\mathrm{hor}}),\label{eq:controlorder}
\end{equation}
their augmented paths satisfy $Z_t\preceq Z'_t$ for all $t\in[0,K_{\mathrm{hor}}]$ almost surely. In particular, one construction simultaneously couples every target set $S\subseteq V$ in Eq.~\ref{eq:control}, and
\begin{equation}
 S\subseteq S'\quad\Longrightarrow\quad
 L(S',\xi)\le L(S,\xi),\qquad Q(S')\le Q(S).
 \label{eq:nestedorder}
\end{equation}
\end{theorem}
\begin{proof}
Increasing $x$ in the product node order decreases each rumor indicator and increases each correction indicator, so $\Phi_\sigma(x)\ge\Phi_\sigma(x')$ whenever $x\preceq x'$. On an interval without node jumps, Eq.~\ref{eq:flow} is a convex combination of the starting memory and the drive. If initially $a\ge a'$ coordinatewise and $x\preceq x'$, it follows that $a(t)\ge a'(t)$ throughout the interval.

At ordered augmented states, every summand in $\ell_i^R$ is nonincreasing: rumor-donor indicators decrease and $1+a_\sigma$ decreases. Every summand in $\ell_i^C$ is nondecreasing: correction-donor indicators increase, $1-a_\sigma$ increases, and the control increases. Nonnegativity allows these factors to be compared termwise. Hence
\begin{equation}
 \ell_i^R(z)\ge\ell_i^R(z'),\qquad
 \ell_i^C(z,u)\le\ell_i^C(z',u').\label{eq:pressureorder}
\end{equation}

Choose finite action-independent bounds $B_i^R$ and $B_i^C$ on the two adoption pressures. For each node and destination $v\in\{R,C\}$, generate an independent rate-$B_i^v$ Poisson proposal process with independent uniform marks. A proposal resets node $i$ to $v$ if its mark is at most $\ell_i^v/B_i^v$. A zero bound has no proposals. Add a common unit-rate proposal process that always resets node $i$ to $U$. All controls use the same proposal times and marks.

At an $R$ proposal, acceptance in the upper path implies acceptance in the lower path, by Eq.~\ref{eq:pressureorder}. The possible outcomes---both nodes reset to the minimal state, only the lower node resets, or neither resets---all preserve order. At a $C$ proposal, acceptance in the lower path implies acceptance in the upper path, and the analogous argument uses the maximal state. A common reset to $U$ is a constant map and preserves order even though it raises an $R$ state and lowers a $C$ state. Memory is continuous at each node jump, so its order is unchanged there. Induction over the almost surely finite proposal sequence, with the flow argument between proposals, proves pathwise comparison.

For each marginal, thinning a rate-$B_i^v$ process with acceptance probability $\ell_i^v/B_i^v$ gives the prescribed conditional adoption intensity. Resets to the current state are inert. Unit-rate resets give both forgetting transitions, and Eq.~\ref{eq:flow} gives the exact deterministic part. Together, the reset intensities and deterministic flow give the specified piecewise-deterministic law, completing the exactness argument.

The collection of all $2^n$ target sets is finite, and the common envelope $u$ bounds all its controls. It can therefore be coupled on one probability space without changing the construction. For nested sets, ordered controls start from the same context. The indicator of state $R$ is order-reversing; integrating it proves Eq.~\ref{eq:nestedorder}, and taking expectations proves the population inequality.
\end{proof}

Although the node process $X$ alone need not be Markov, $(X,A)$ determines both the memory flow and the jump intensities. Conditioning on a history therefore supplies the augmented initial state needed for the pathwise argument. The construction also compares distinct histories when their augmented states are ordered; only the queried controls need to be simulated.

\begin{corollary}[Heterogeneous relaxation and finite higher-order factors]
\label{cor:general}
The construction extends to a finite family of memory coordinates $a_f\in[-T_{0,f},T_{0,f}]$ with common-across-actions rates $\lambda_f\ge0$ and dynamics
\begin{equation}
 \dot a_f=\lambda_f[\Phi_f(x)-a_f],\qquad
 |\Phi_f(x)|\le T_{0,f},\label{eq:generalflow}
\end{equation}
provided each drive $\Phi_f$ is antitone in the node order, rumor reset intensities are antitone in the augmented order, correction reset intensities are monotone, and the intensities admit common finite bounds. Common deterministic, bounded, time-dependent reset-to-$U$ rates are also allowed. In particular, the result covers nonnegative weighted edges and finite hyperedges with all-rumor/all-correction donor products and memory drives
\begin{equation}
 \Phi_f(x)=T_{0,f}\big[\ind\{N_R(f,x)\ge q_f^R\}
              -\ind\{N_C(f,x)\ge q_f^C\}\big],\label{eq:hyperdrive}
\end{equation}
where $0<T_{0,f}\le1$ and the fixed thresholds $q_f^R,q_f^C$ are positive integers.
\end{corollary}
\begin{proof}
For each coordinate, the flow is Eq.~\ref{eq:flow} with its own $\lambda_f$ and drive. Both coefficients remain nonnegative, so unequal relaxation rates across coordinates preserve the same memory order. The event argument in Theorem~\ref{thm:coupling} uses only the two intensity inequalities and a common constant reset map. A bounded deterministic forgetting rate can be generated by its own common time-inhomogeneous Poisson process. Common finite bounds give finitely many proposals on the finite horizon.

For Eq.~\ref{eq:hyperdrive}, increasing node states decreases $N_R$ and increases $N_C$, so the drive is antitone and lies in $[-T_{0,f},T_{0,f}]$. Nonnegative weighted sums of donor products have the required intensity order. The factors $1\pm a_f$ are nonnegative because $T_{0,f}\le1$. These properties verify every assumption for the stated higher-order family.
\end{proof}

The essential conditions are the order of the memory drive and reset intensities. Triangles provide the concrete model tested below; finite weighted hyperedges satisfy the same comparison argument when these conditions hold.

\subsection{Fluctuations of nested response differences}
\begin{theorem}[Range--mean variance and sufficient sampling bound]
\label{thm:variance}
For fixed nested sets $S\subseteq S'$, let
$Y=L(S,\xi)-L(S',\xi)$ under Theorem~\ref{thm:coupling} and let
$\mu=Q(S)-Q(S')$. Then $Y\in[0,H]$ and
\begin{equation}
 \Var(Y)\le\mu(H-\mu).\label{eq:variance}
\end{equation}
If $\mu=0$, then $Y=0$ almost surely. If $\mu>0$, $0<\varepsilon<1$, and $Y_1,\ldots,Y_m$ use independent tapes, their mean $\bar Y_m$ satisfies
\begin{equation}
 \Prob\{ |\bar Y_m-\mu|\ge\varepsilon\mu\}
 \le2\exp\!\left[-\frac{m\varepsilon^2\mu}
 {2(H-\mu)+(2/3)H\varepsilon}\right].\label{eq:relative}
\end{equation}
A failure probability at most $\delta\in(0,1)$ is ensured by
\begin{equation}
 m\ge\frac{2(H-\mu)+(2/3)H\varepsilon}{\mu\varepsilon^2}
       \log\frac2\delta.
 \label{eq:samples}
\end{equation}
A simpler sufficient choice is $m\ge[8H/(3\mu\varepsilon^2)]\log(2/\delta)$.
\end{theorem}
\begin{proof}
Pathwise comparison gives $0\le Y\le H$, so $Y^2\le HY$; taking expectations and subtracting $\mu^2$ proves Eq.~\ref{eq:variance}. The zero-mean case follows because a nonnegative random variable with zero expectation vanishes almost surely. Bernstein's bounded-variable inequality, with variance at most $\mu(H-\mu)$ and centered range bounded by $H$, gives
\[
 \Prob\{|\bar Y_m-\mu|\ge t\}
 \le2\exp\!\left[-\frac{mt^2}{2\mu(H-\mu)+(2/3)Ht}\right].
\]
Substituting $t=\varepsilon\mu$ gives the stated condition on $m$, with the simpler bound following from
$2(H-\mu)+(2/3)H\varepsilon\le8H/3$.
\end{proof}

Path order makes the classical range--mean inequality applicable to intervention contrasts. It is sharp on $[0,H]$: $Y=H Z_{\mathrm{Ber}}$, with $Z_{\mathrm{Ber}}$ Bernoulli of mean $\mu/H$, attains equality. This establishes sharpness over bounded random variables, not attainability within this particular spreading model. At fixed confidence, the sufficient sample scale is $H/(\mu\varepsilon^2)$ for the ordered contrast. For a generic signed contrast in $[-H,H]$, the worst-case variance bound $H^2$ instead gives $H^2/(\mu^2\varepsilon^2)$. At fixed $H$ and $\varepsilon$, dependence on $1/\mu$ is linear in the ordered case and quadratic in the generic case. These are sufficient upper bounds, not matched upper and lower complexity bounds or lower bounds for RTC. Since Eq.~\ref{eq:samples} involves the unknown mean, using it as a stopping rule requires further statistical control.

\begin{corollary}[Other monotone occupation costs]
Let $c_t(z)$ be a bounded order-reversing running cost, and suppose its integrated loss lies in $[0,M]$. For ordered controls, the coupled integrated losses are ordered and Theorem~\ref{thm:variance} holds with $M$ in place of $H$.
\end{corollary}
\begin{proof}
The path order gives $c_t(Z'_t)\le c_t(Z_t)$ at every time. Integrate and apply the proof of Theorem~\ref{thm:variance} to the resulting difference in $[0,M]$.
\end{proof}

\subsection{Ordered envelopes for equal-cardinality comparisons}
Two distinct size-$k$ actions cannot be nested. Their intersection and union nevertheless provide ordered auxiliary controls that bound both losses on each tape.
\begin{theorem}[A variance bound for nonnested actions]
\label{thm:exchange_envelope}
Fix an initial augmented state and any two target sets $S,T_{\mathrm{set}}$. On the grand graphical coupling, put $I=S\cap T_{\mathrm{set}}$, $J=S\cup T_{\mathrm{set}}$,
\begin{align*}
 W&=L(S)-L(T_{\mathrm{set}}), &B_{\mathrm{env}}&=L(I)-L(J),\\
 P_S&=L(S)-L(J), &P_{T_{\mathrm{set}}}&=L(T_{\mathrm{set}})-L(J).
\end{align*}
Write $\Delta_{SW}=\E W$, $b=\E B_{\mathrm{env}}$, and $v=\E(P_S+P_{T_{\mathrm{set}}})$. Then
\begin{equation}
 |W|\le B_{\mathrm{env}}\le H,\qquad
 \Var(W)\le H\min\{b,v\}-\Delta_{SW}^2.
 \label{eq:exchange_var}
\end{equation}
In particular, no nesting of $S,T_{\mathrm{set}}$ is required. If $b=0$, their coupled losses agree almost surely. For independent tapes and $\Delta_{SW}\ne0$, a sufficient condition for relative error $\varepsilon|\Delta_{SW}|$, at confidence $1-\delta$, with $\varepsilon>0$ and $0<\delta<1$, is
\begin{equation}
 N\ge \frac{2H\min\{b,v\}+\frac43H\varepsilon|\Delta_{SW}|}
 {\varepsilon^2\Delta_{SW}^2}\log\frac2\delta.
 \label{eq:exchange_samples}
\end{equation}
The inequality assumes the same per-target intensity and duration for all four sets; the union is an auxiliary simulation, not a feasible size-$k$ deployment.
\end{theorem}
\begin{proof}
Nested-set comparison gives $L(J)\le L(S),L(T_{\mathrm{set}})\le L(I)$ on each tape. Two numbers in this random interval differ by at most its length, proving $|W|\le B_{\mathrm{env}}\le H$. Thus $W^2\le B_{\mathrm{env}}^2\le HB$. Also $P_S,P_{T_{\mathrm{set}}}\ge0$, $P_S,P_{T_{\mathrm{set}}}\le H$, and $W=P_S-P_{T_{\mathrm{set}}}$, so
$W^2=P_S^2+P_{T_{\mathrm{set}}}^2-2P_SP_{T_{\mathrm{set}}}\le H(P_S+P_{T_{\mathrm{set}}})$. Taking both expectations and subtracting $\Delta_{SW}^2$ proves Eq.~\ref{eq:exchange_var}. A nonnegative random variable with mean zero vanishes almost surely, giving the zero-envelope claim. Finally $W\in[-H,H]$ implies $|W-\Delta_{SW}|\le2H$. Bernstein's inequality, using $\Var(W)\le H\min\{b,v\}$ and deviation $\varepsilon|\Delta_{SW}|$, yields Eq.~\ref{eq:exchange_samples}.
\end{proof}
For equal-cardinality exchanges, the bound is useful when the expected auxiliary loss interval is narrow. Its quantities $b,v,\Delta_{SW}$ are population values, and estimating the intersection and union incurs additional simulation cost.

\subsection{Strict complementarity and the limit of submodular guarantees}
For a monotone set function $g$, define its truncated submodularity ratio by
\begin{equation}
 \gamma_g=\min\!\left\{1,\inf_{\substack{L\cap A=\varnothing\\g(L\cup A)>g(L)}}
 \frac{\sum_{v\in A}[g(L\cup\{v\})-g(L)]}{g(L\cup A)-g(L)}\right\}.
 \label{eq:ratio}
\end{equation}
Restricting the infimum to sets with the desired cardinality bounds gives the corresponding restricted ratio. A positive lower bound is the hypothesis used by weak-submodularity greedy guarantees~\cite{das2018}.

\begin{proposition}[A complete-process zero-ratio example]
\label{prop:counter}
On one filled triangle $\{i,j,v\}$, initialize $(X_i,X_j,X_v)=(U,U,R)$ and $A_0=0$. Let $\beta_1=0$, $\beta_2,u,D,H>0$, $0<T_0\le1$, and $\lambda\ge0$. Then
\begin{equation}
 g(\{i\})=g(\{j\})=0,\qquad g(\{i,j\})>0.
 \label{eq:counter}
\end{equation}
Consequently $\Gamma_{ij}(\varnothing)>0$ and $\gamma_g=0$.
\end{proposition}
\begin{proof}
Initially a single rumor node is present, and the two-donor requirement prevents further rumor transmissions. With no intervention or either singleton intervention, two correction donors cannot form while $v$ remains in $R$. Because its only exit from $R$ is the independent unit-rate forgetting event, all three expected losses are
\begin{equation}
 Q(\varnothing)=Q(\{i\})=Q(\{j\})
       =\frac{e^{-D}(1-e^{-H})}{3}.\label{eq:singletonloss}
\end{equation}

With both $i$ and $j$ targeted, there is positive probability that both adopt $C$ before $D$, stay in that state over a nonempty interval, and $v$ has not yet forgotten. No path has two rumor nodes, so its memory drive is always nonpositive. Since $A_0=0$, Eq.~\ref{eq:flow} gives $A(t)\le0$. During the simultaneous-correction interval, $v$ consequently has additional correction rate at least $\beta_2T_0>0$. There is positive probability that this reset precedes natural forgetting, that natural forgetting would occur after $D$, and that a positive portion of the post-control rumor occupation is removed. This event has positive probability for every $D,H>0$ and every $\lambda\ge0$, including $T_0=1$.

Theorem~\ref{thm:coupling} rules out compensating increases on other tapes. The pair thus has strictly smaller expected loss than no intervention. Equation~\ref{eq:counter} follows, and the ratio in Eq.~\ref{eq:ratio} is zero at $L=\varnothing$, $A=\{i,j\}$.
\end{proof}

\begin{corollary}[No uniform ratio over positive pairwise rates]
\label{cor:positivebeta}
Keeping the other parameters of Proposition~\ref{prop:counter} fixed, the submodularity ratios can approach zero as strictly positive $\beta_1\downarrow0$. Hence no uniform positive ratio applies to the entire positive-pairwise-rate model class.
\end{corollary}
\begin{proof}
Let $m_E$ be the number of undirected edges and $K_{\mathrm{tot}}=D+H$. Couple the zero-rate process and a process with pairwise rate $b>0$ by adding donor arrows of rate $bT$ for each direction and each of the two transmission types on every edge. The total added-arrow rate is $4m_EbT$. If no added arrow occurs by $K_{\mathrm{tot}}$, the two processes, including memory, coincide. For every action $S$,
\begin{equation}
 |Q_b(S)-Q_0(S)|\le H p_b,\qquad
 p_b=1-e^{-4m_EbTK}\longrightarrow0.\label{eq:continuity}
\end{equation}
Let $c=g_0(\{i,j\})>0$. Each singleton benefit is at most $2Hp_b$, and the joint benefit is at least $c-2Hp_b$. For small enough $b$, Eq.~\ref{eq:ratio} is bounded above by
$4Hp_b/(c-2Hp_b)$, which tends to zero.
\end{proof}

The two targeted nodes jointly open a correction channel that neither can open alone. The absence of pairwise transmission and the initial donor configuration make the singleton gains vanish; positive $\beta_2,u,D,H$ ensure a positive joint gain. The example separates monotone response from diminishing returns: individually ineffective perturbations can jointly create the donors needed to change a third node. The zero ratio prevents a direct weak-submodularity guarantee for the full model class. Approximation results therefore require additional structure, as in the restricted score analysis developed in the supplement.

\subsection{Memory generated by the intervention}
Couple an untreated path $Z^0$ and a path $Z^S$ controlled on $[0,D)$ from the same initial state. At withdrawal, let
\[
 Z^{h}_D=(X^S_D,A^0_D).
\]
The hybrid keeps the treated node configuration and replaces only its memory by the paired untreated memory. Theorem~\ref{thm:coupling} gives
\[
 Z^0_D\preceq Z^{h}_D\preceq Z^S_D.
\]
Continue the three states without control, using the same fresh post-withdrawal tape. The theorem applies again, so
\begin{equation}
 L^S\le L^{h}\le L^0,\qquad
 L^0-L^S=(L^0-L^{h})+(L^{h}-L^S).\label{eq:withdrawal_decomposition}
\end{equation}
The first term measures the benefit carried by the treated node state at baseline memory; the second measures the extra benefit of retaining the intervention's memory. Both terms are nonnegative under this coupling. Write $G_X=\E(L^0-L^{h})$, $G_A=\E(L^{h}-L^S)$ and $G=G_X+G_A$. For $G>0$, define
\begin{equation}
 \eta_A=\frac{G_A}{G}\in[0,1].\label{eq:memory_fraction}
\end{equation}
If $G=0$, the fraction is undefined and is not assigned zero. This is a model-internal, paired counterfactual decomposition: the baseline memory is taken from its coupled untreated history, rather than estimated from an observed social intervention. Strict positivity is a numerical question, not a consequence of monotonicity alone. The hybrid combines variables from two coupled histories, so its expectation and the fraction refer to the specified common-event pairing; they are not a coupling-invariant causal mediation estimand.

\section{Numerical construction and observables}
\label{sec:simulation}
\subsection{Exact reset construction}
Let $d_i$ and $f_i$ denote the numbers of edges and filled triangles incident to node $i$. For the intervention family in Eq.~\ref{eq:control}, use
\begin{align}
 B_i^R&=\beta_1Td_i+\beta_2T_0(1+T_0)f_i,\notag\\
 B_i^C&=B_i^R+u,\qquad B_i^U=1.\label{eq:bounds}
\end{align}
These bounds are shared by all actions for the whole horizon. The correction envelope includes $u$ even at untargeted nodes and after withdrawal. This keeps proposal times independent of the deployed action. Proposal channels are indexed by node and destination. Counter-indexed random variates specify each channel's exponential increments and acceptance marks; different actions can therefore replay the same tape without storing it or changing its alignment when acceptance decisions differ.



For each face the implementation stores its last advanced memory, that time, and its current drive. The drive changes only when one of its three nodes changes. A pressure query evaluates the memory analytically at the current time; an accepted state change advances and refreshes only incident faces. These analytic updates advance memory exactly between node transitions.

For each marginal path, proposals are accepted against its current intensity and memory is advanced by the exact exponential flow. Self resets leave the state unchanged. A complete state trajectory continues after rumor extinction to record memory; a loss-only continuation may stop there because no rumor donor remains. The implementation and all pseudocode are in Supplementary Section~\ref{S-sec:full_computational_application}.

\section{Physical response after withdrawal}\label{sec:experiments}
\subsection{Design and numerical checks}
The new withdrawal study uses eight synthetic graphs ($n=100$, two per Erd\H{o}s--R\'enyi, Watts--Strogatz, sparse blockmodel and preferential-attachment-with-closure family) and two observed contact topologies. Each context starts with $\lfloor0.1n\rfloor$ rumor nodes, the same number of correction nodes, and zero memory; two uncontrolled time units generate a history. Each graph or topology supplies two histories. All histories, including extinction, are used. Correction targets four nodes, ranked first by current rumor status and then by degree, with deterministic ties. We use $T_0=0.9$, $\lambda=0.2$, $u=3$, $H=10$ and $D\in\{1,2,4\}$. Rates satisfy $\beta_1T_0\bar d=1.15$ and $\beta_2T_0\bar f=3$, where $\bar d$ and $\bar f$ are mean edge degree and face participation.

The synthetic contexts use 64 paired replications per duration and history; each observed topology uses 128. Untreated and treated paths share the pre-withdrawal events. After withdrawal, the untreated, retained-memory, baseline-memory and zero-memory branches share a fresh event tape. The retained and baseline-memory treated branches have exactly the same node state at withdrawal. Configurations, source hashes and all reporting rules were fixed before the new runs; these studies are descriptive and are separate from earlier confirmation families. Synthetic uncertainty resamples graphs within families; observed-topology intervals resample paired replications within each fixed history, with 2,000 bootstrap repetitions. Pointwise intervals do not represent uncertainty over a population of contact networks.

Exact checks link the implementation to the mathematical comparisons. The arity-general simulator agrees with the archived triangle simulator on 256 loss samples, with zero discrepancy. Earlier 27-state frozen-memory triangle calculations agree by matrix exponentials and uniformization to $3.61\times10^{-16}$. The withdrawal trajectories satisfy Eq.~\ref{eq:withdrawal_decomposition} on every sampled path. Full parameters, topology processing and implementation checks are in Supplementary Section~\ref{S-sec:physics_new}.

\subsection{The memory contribution generated by correction}
Figure~\ref{fig:mediation} separates the total occupation-loss reduction into the node-state and memory components of Eq.~\ref{eq:withdrawal_decomposition}. Across the eight synthetic graphs, mean memory contributions are $0.01986$, $0.03634$ and $0.06849$ at $D=1,2,4$, respectively. Their shares of total gain are $3.51\%$, $4.51\%$ and $6.64\%$. Memory contributions are positive on eight, seven and seven graphs, respectively; the remaining graph-level means are zero. The duration dependence reported here concerns the tested grid.

The hybrid branch supplies the relevant counterfactual. It leaves the treated node configuration unchanged and replaces memory by the paired untreated value. Its higher loss therefore quantifies a benefit carried by the intervention's memory after the external input has ended. It does not require the retained memory itself to be negative or beneficial relative to erasing all memory. Table~\ref{tab:mediation} reports the $D=2$ comparison, while every duration and context is given in the supplement.

\begin{figure*}[t]\centering\includegraphics[width=.98\textwidth]{physicaA_memory_mediation.pdf}
\caption{Memory generated by a temporary intervention. (a) Mean total occupation-loss reduction on eight synthetic graphs, decomposed into node-state and memory components. (b) Memory-mediated fractions for the specified paired protocol on two observed contact topologies, reported separately. Error bars are descriptive pointwise 95\% bootstrap intervals conditional on each topology and its two histories. Lines connect tested durations. Rates are synthetic in all panels.}\label{fig:mediation}\end{figure*}
\begin{table*}[t]\centering\small
\caption{Withdrawal-time decomposition at $D=2$. Loss has time units. Intervals are descriptive pointwise 95\% intervals: graph-stratified for synthetic networks and conditional Monte Carlo intervals for each observed topology. The memory fraction is a ratio of mean gains.}\label{tab:mediation}
\resizebox{\linewidth}{!}{\begin{tabular}{lrrrrr}\toprule
Topology & Nodes & Total gain & Node-state gain & Memory gain & Memory fraction (\%)\\\midrule
Synthetic, eight graphs & 100 & 0.80570 & 0.76936 & 0.03634 & 4.51 [3.11, 5.84]\\
High-school contacts & 327 & 0.12825 & 0.10364 & 0.02462 & 19.19 [11.39, 28.94]\\
Primary-school contacts & 242 & 0.80637 & 0.75449 & 0.05187 & 6.43 [2.77, 12.90]\\\bottomrule
\end{tabular}}\end{table*}

\subsection{Observed contact topologies}
The contact groups come from the public simplicial-closure collection~\cite{benson2018}, based on the high-school and primary-school sensor studies~\cite{highschool2015,primaryschool2011}. In that collection, groups are maximal cliques within 20-second contact windows. We use all 327 high-school participants and all 242 primary-school participants. Pairwise edges are projected from the recorded groups; filled faces are unique groups recorded with exactly three members. We do not fill every triangle of the time-aggregated graph. The resulting complexes have 5,818 edges and 2,091 faces for the high school, and 8,317 edges and 4,600 faces for the primary school.

Figure~\ref{fig:contacts} shows the post-withdrawal response for $D=2$. The two treated branches start from an identical node configuration, yet retaining intervention-generated memory gives lower residual loss than replacing it by baseline memory. The memory-mediated loss differences are $0.02462$ $[0.01382,0.03889]$ and $0.05187$ $[0.02191,0.10164]$ for the high-school and primary-school topologies. Their fractions differ substantially and are reported separately. All durations and zero-memory controls are reported separately in Supplementary Section~\ref{S-sec:physics_new}.

\begin{figure*}[t]\centering\includegraphics[width=.98\textwidth]{physicaA_contacts.pdf}
\caption{Same treated node state, different memory at withdrawal. Curves show the two observed contact topologies at $D=2$ under synthetic rates. Retained-memory and baseline-memory treated branches coincide at time zero because their node states are identical; their subsequent responses differ. The untreated branch also differs in its withdrawal node state. Each curve averages 128 paired replications for each of two fixed histories.}\label{fig:contacts}\end{figure*}

\subsection{Forgetting an initial bias versus following subsequent activity}
Earlier controlled-memory comparisons hold higher-order transmission active and distinguish $A_\sigma\equiv0$, frozen memory and dynamic memory at $\lambda=0.02,0.2,2$. Freezing means $\lambda=0$ and removes modulation only when the initial memory is also zero. The complete curves and all settings are in Supplementary Section~\ref{S-sec:earlier_memory_cohorts}.

A separate 32-graph natural-history study finds that the controlled mean loss increases from $1.63886$ to $2.04232$ as the continuation relaxation rate increases from $0.02$ to $2$, while the loss difference between natural and reset dynamic initial memory falls from $0.19598$ to $0.02818$. This separates sensitivity to the initial memory from the burden sustained by subsequent feedback. At $\lambda=2$ and the observation endpoint, the mean direct initial-memory term is approximately $1.69\times10^{-13}$, but the controlled feedback residual is $0.00907$ $[0.00477,0.01387]$. The corresponding uncontrolled interval includes zero. These descriptive results and the full time grid are reported in Supplementary Section~\ref{S-sec:natural_memory_details}.

\subsection{A finite-time retention and rebound map}
On one triangle we initialize $(R,R,U)$ and positive memory $T_0$, apply correction to the first two nodes, and scan the complete grid $\lambda\in\{0.02,0.05,0.1,0.2,0.5,1,2,5\}$ and $D\in\{0.25,0.5,1,2,4,8\}$. Each cell has 512 paired replications and a post-withdrawal sampling interval of 0.05. We define pathwise sampled rebound as
\begin{equation}
 \mathcal R=\E\!\left[\max_{t\in\mathcal T_H}\rho_R(D+t)-\rho_R(D)\right],\qquad
 \mathcal T_H=\{0,0.05,\ldots,10\}.\label{eq:rebound}
\end{equation}
Extinction trajectories contribute zero and are not removed. The maximum of the mean occupation curve is a different observable and is reported separately.

Figure~\ref{fig:triangle} contrasts direct-memory retention $e^{-\lambda D}$ with the measured response. The direct term changes smoothly around $\lambda D$ of order one. Sampled pathwise rebound ranges from zero to $0.13997$ and is mainly duration-dependent on this grid. The overshoot of the mean curve is at most $0.00846$, and the memory-mediated fraction is at most $1.96\%$. Thus individual-path rebound is not evidence for an ensemble rebound or a critical boundary. The small triangle illustrates the retention scale and finite-state fluctuations; the larger-network withdrawal experiment supplies the stronger memory-mediation evidence.

\begin{figure*}[t]\centering\includegraphics[width=.98\textwidth]{physicaA_triangle_crossover.pdf}
\caption{One-triangle finite-time response. (a) Analytic direct-memory retention $e^{-\lambda D}$. (b) Measured mean pathwise sampled rebound, Eq.~\ref{eq:rebound}. (c) Memory-mediated gain fraction, expressed in percent. Each panel uses its stated linear color scale over the full tested grid. The dashed line $\lambda D=1$ is a time-scale guide, not an inferred transition boundary.}\label{fig:triangle}\end{figure*}

\subsection{Comparison bounds and higher-order factors}
Across 72 earlier graph--memory contexts, the median empirical variance divided by the plug-in bound is $0.16093$ for nested contrasts and $0.11129$ for equal-size exchanges sharing three targets. All denominators exceed $10^{-12}$. These ratios describe slack in conservative bounds; because their moments use the same samples, inequalities for the empirical distribution follow algebraically and are not independent tests of population coverage. Full diagnostics are in Supplementary Section~\ref{S-sec:response_extension}.

A weighted 12-node network with four-uniform hyperedges provides a small check beyond triangles. With nonnegative weights, the all-donor reset rates and either a majority drive or the bounded linear drive $\Phi_f=T_{0,f}(N_R-N_C)/|f|$ satisfy Corollary~\ref{cor:general}. On 256 shared replications per control and drive, all nested occupation inequalities, terminal node and reverse-memory orders, and nonnested intersection--union envelopes hold. These examples illustrate the sufficient conditions in a setting beyond triangles. The complete setup and curves are in Supplementary Fig.~\ref{S-fig:hyperedge4}.

\section{Joint evaluation as a computational application}
The comparison law also permits exact shared evaluation of candidate controls. The following statement concerns the joint process, not only its marginals.
For a finite portfolio of $M$ actions write $z=(z^1,\ldots,z^M)$. A reset channel $c=(i,d)$ has effective intensity
$q_c^j(z^j,t)=\ell_i^d(z^j,t)\ind\{x_i^j\ne d\}$, with $\ell_i^U=1$.
Let $F_c(z,h)$ reset copy $j$ to destination $d$ exactly when $h<q_c^j(z^j,t)$. Self resets are inert. These conventions define a common marked-Poisson representation of the type used in stacked coupling.
\begin{theorem}[Exact joint simulation with predictable bounds]
\label{thm:batch}
Assume finitely many channels, bounded intensities, the deterministic memory flow in Eq.~\ref{eq:flow}, and a locally finite refresh schedule. Require finite predictable proposal bounds whose total intensity is integrable on every finite time interval, so auxiliary proposals do not accumulate. Between refresh times choose predictable constants $b_c\ge\max_j q_c^j$ that remain valid until the next refresh or scheduled control boundary. Simulate proposals at total rate $\sum_cb_c$, choose channel $c$ with probability $b_c/\sum_rb_r$, and apply $F_c(z,h)$ with a uniform mark $h\in[0,b_c]$. Refresh affected bounds after jumps and handle control boundaries explicitly. This algorithm has the same joint path law as the fixed-envelope graphical construction on the queried portfolio. Deleting channels that are inert in every copy does not change this law. Restricting the portfolio to any subset gives the same joint law as evaluating that subset alone.

An adaptive finite-round search has the same transcript and returned-action distribution under either evaluator, provided each round fixes its portfolio from the preceding transcript and uses fresh independent replication tapes. This is distributional equality; different evaluators need not return the same realization for an identical integer seed.
\end{theorem}
\begin{proof}
For a smooth test function $f$, thinning produces the jump term
\begin{align*}
 &\sum_{c:b_c>0}\int_0^{b_c}[f(F_c(z,h))-f(z)]\,dh\\
 &\qquad=\sum_c\int_0^\infty[f(F_c(z,h))-f(z)]\,dh.
\end{align*}
The equality follows because $F_c(z,h)=z$ above the maximum effective intensity. Adding the deterministic memory drift gives exactly the generator of the fixed marked-Poisson construction. The bounds cannot be crossed before refresh by assumption; scheduled discontinuities are treated as deterministic boundaries. The integrable proposal envelopes and locally finite refreshes exclude accumulation of auxiliary events. Bounded physical rates exclude accumulation of state jumps. Deterministic flow with these bounded jump kernels gives a unique recursively constructed joint path law. A jointly inert channel has identically zero integrand. For $f$ depending on only a subset of copies, the integral depends only on that subset's intensities and reset states; hence its projected process has the same unique law. Conditioning on the preceding transcript fixes the next portfolio. Fresh replication tapes then give identical conditional observation laws under the two evaluators, so induction over rounds completes the adaptive argument.
\end{proof}
The supplementary search example uses a frozen-memory score that treats targets as corrected donors and combines direct, edge and face correction exposure through a saturating function. It proposes candidates; their losses are evaluated with the full dynamic process. The implementation shares event scheduling without changing the joint observation law supplied to a selection procedure. In a 32-graph whole-search study, its geometric runtime was 73.6\% lower than persistent-topology action-wise replay. That comparison includes amortized preparation and interface costs. It is distinct from comparing couplings with different joint laws; full timing and search results, including inconclusive and unfavorable comparisons, are in Supplementary Section~\ref{S-sec:full_computational_application}.

At the earlier four-action quality target of mean excess loss at most $0.02$, the ordered construction used 128 physical trajectories versus 512 for global event thinning, a 75\% reduction on the discrete budget grid. Common reaction paths (RTC) reached the target at 128 trajectories as well. The crossing uses a mean point estimate relative to a finite independent reference, not a global optimum or a per-query confidence certificate. The threshold-free extension below avoids treating that crossing as a general efficiency law.

\subsection{Precision without a quality threshold}
Four fixed graph/history contexts, one per network family, compare two distinct size-four actions. For each method, eight selection-free repetitions use 8, 16, 32, 64 and 128 paired samples. An independent 2,048-pair reference estimates the expected contrast. We report squared error relative to this finite reference, paired-loss variance divided by sample count, and measured computation cost. The reference has Monte Carlo uncertainty; its squared error is not exact oracle MSE. All budgets and contexts are reported in Fig.~\ref{fig:precision} and Supplementary Section~\ref{S-sec:precision_new}.

At 128 pairs, the context-averaged reference-based MSE is 0.01491 for thinning, 0.00269 for RTC and 0.00284 for the ordered construction. Corresponding paired-loss sample variances are 2.21746, 0.40761 and 0.31709. The variance reduction relative to thinning therefore extends beyond one quality threshold; the small MSE difference from RTC does not establish a precision advantage over RTC. Runtime comparisons are descriptive implementation measurements.
\begin{figure*}[t]\centering\includegraphics[width=.98\textwidth]{physicaA_precision.pdf}
\caption{Threshold-free contrast estimation on four fixed graph/history contexts. (a) Mean squared error relative to an independent finite reference, averaged across contexts and eight repetitions. (b) Sample paired-loss variance divided by pair budget. (c) The same variance scale against measured wall time. Both actions count toward physical trajectory cost. Reference uncertainty and all graph-level values are reported in the supplementary data.}\label{fig:precision}\end{figure*}

\section{Discussion}
The withdrawal comparison identifies a memory contribution without changing the treated node state. This resolves a distinction that prescribed $\pm T_0$ initializations cannot: the extra response is carried by memory generated along the intervention, rather than by an arbitrary initial bias. Positive contributions occur in the tested synthetic and contact-topology settings, with larger shares at longer durations on the observed grid. Their magnitudes remain topology- and protocol-dependent.

Memory relative to an untreated history and memory relative to complete erasure are different comparisons. On both contact topologies, setting memory to zero can yield smaller loss than retaining treated memory, even though retaining treated memory improves on resetting it to the untreated level. An intervention can therefore leave memory more favorable than the baseline while residual memory remains unfavorable compared with erasure. The mechanism concerns the direction and persistence of group influence, not an assertion that memory always supports correction.

Along a realized node path,
\[
 A_\sigma(t)=e^{-\lambda t}A_\sigma(0)
 +\lambda\int_0^t e^{-\lambda(t-s)}\Phi_\sigma(X(s))\,ds.
\]
The first term gives direct forgetting; the second stores subsequent group activity. Faster relaxation reduces dependence on the first term while strengthening tracking of the second. The natural-history and withdrawal experiments consequently ask different questions. The one-triangle map separates the analytic retention scale from sampled rebound, and supplies no evidence for a thermodynamic phase transition or a critical point.

Theorems~\ref{thm:coupling} and~\ref{thm:exchange_envelope} give sufficient comparison conditions and conservative bounds. We have not proved their necessity. Weighted four-uniform examples illustrate the extension in Corollary~\ref{cor:general}; other bounded antitone drives, including a linear count drive, fit the same order argument. State-dependent forgetting, control-dependent relaxation or redistribution of a fixed total control intensity would require a new compatibility analysis.

The model is a theoretical abstraction. Transition rates and memory scales are prescribed, and sensor-derived groups provide topology rather than observations of misinformation or correction. The empirical-topology results are conditional on two histories per dataset and are reported separately. Finite-time occupation, pathwise rebound and a memory-mediated fraction are distinct observables. Approximation guarantees for intervention search remain a separate problem: the counterexample rules out a uniform positive submodularity ratio, and the supplementary search results have mixed transfer beyond their calibrated settings.

\section{Conclusion}
A temporary intervention can leave a response in group memory as well as in node states. An ordered withdrawal-time comparison separates these contributions, and simulations on synthetic networks and two observed contact topologies quantify a nonzero memory-mediated component in the tested settings. The augmented node--memory order yields exact stochastic comparisons, conservative fluctuation bounds and nonnested envelopes, while a complete-process counterexample separates monotonicity from diminishing returns. Direct memory retention and finite-time rebound have different time-scale dependence, so faster forgetting of initial memory does not by itself imply a smaller spreading burden. Exact joint evaluation supports the comparisons computationally; optimization and runtime results remain secondary applications with their stated limits.

\section*{CRediT authorship contribution statement}
\textbf{Junkai Fan:} Conceptualization, Methodology, Software, Formal analysis, Investigation, Validation, Data curation, Visualization, Writing -- original draft, Writing -- review \& editing.
\section*{Acknowledgements and funding}
This research received no external funding.
\section*{Declaration of competing interest}
The author declares no competing interests.
\section*{Data and code availability}
The reproducibility package for this revision is assigned Zenodo DOI \href{https://doi.org/10.5281/zenodo.23811459}{10.5281/zenodo.23811459}. It contains raw simulation outputs, configuration files, seeds, analysis scripts and numerical source maps. Versioned code is available in the \path{physicaA-submission-20261031} branch of \href{https://github.com/Lucas-fan666/code/tree/physicaA-submission-20261031}{Lucas-fan666/code}. The archive manifest identifies the original study files and the new physical-response experiments. Contact data are credited to the simplicial-closure collection and original sensor studies~\cite{benson2018,highschool2015,primaryschool2011}.
\section*{Declaration of generative AI and AI-assisted technologies}
During preparation, the author used ChatGPT (OpenAI) to assist with language revision, LaTeX preparation, research-code development and checking, and programmatic reconstruction of the explanatory illustration. Numerical results were obtained by executing the archived simulation and analysis code. The author checked the core mathematical proofs, reviewed the manuscript and takes full responsibility for its content.
\begingroup
\interlinepenalty=10000
\renewcommand{\bibfont}{\footnotesize}\setlength{\bibsep}{0pt}
\bibliographystyle{elsarticle-num}
\bibliography{references}
\endgroup
\end{document}
